\chapter{Quantisation}\label{ch:quantisation}

Four-bit weights take a decode step from 26.93~GB to 15.01~GB (\arith), and on this board bytes are time. This chapter covers the formats the engine reads, the quantiser that wrote the published 4-bit set, the gate that decides whether a weight set ships, and what stays at full precision.

The served weights are the published set, \code{0xBakeer/TandemLLM-Qwen3.8-27B-NVFP4} on Hugging Face: every projection in NVFP4, quantised once from Qwen's BF16 release with GPTQ, plus an output head in FP8 (e4m3). The files load over the vendor's FP8 checkpoint, which supplies every tensor the overlays do not replace.

\section{What the checkpoint stores}

Qwen ships the model as block-scaled FP8\index{FP8}. Every large projection is a pair of tensors: \code{X.weight} holds e4m3 codes of shape $[N, K]$, and \code{X.weight_scale_inv} holds one BF16 scale per 128 by 128 block. A weight is
\[
w = \textit{code} \cdot \textit{scale}[\lfloor n/128 \rfloor, \lfloor k/128 \rfloor].
\]
\code{engine/loader.py} keeps that pair as an \code{FP8Block} and never expands it in memory, because a decode step pays for exactly these bytes. Some tensors stay BF16 in the checkpoint: the output head, the embedding, every norm, the convolution weights and two small projections of the recurrent layers.

\section{NVFP4}\index{NVFP4}

NVFP4 stores a weight in three tensors:

\begin{table}[htbp]
\centering\small
\caption{The NVFP4 layout of one projection.}\label{tab:nvfp4-layout}
\begin{tabularx}{\linewidth}{@{}L{.22\linewidth}L{.1\linewidth}L{.12\linewidth}Y@{}}
\toprule
\hd{tensor} & \hd{type} & \hd{shape} & \hd{meaning} \\
\midrule
\code{weight} & uint8 & $[N, K/2]$ & two e2m1 codes a byte, low nibble first \\
\code{weight_scale} & e4m3 & $[N, K/16]$ & one scale per 16 weights of a row \\
\code{weight_scale_2} & fp32 & one number & the scale of the scales \\
\bottomrule
\end{tabularx}
\end{table}

A weight is $\text{e2m1}(\textit{code}) \cdot \textit{weight\_scale}[n, \lfloor k/16 \rfloor] \cdot \textit{weight\_scale\_2}$. The e2m1 grid has eight magnitudes (0, 0.5, 1, 1.5, 2, 3, 4 and 6) and a sign. A weight costs 0.5625 bytes against 1.0 in FP8, so one layer's MLP drops from 267.4~MB to 150.4~MB (\arith).

The group of 16 is what makes 4 bits usable. An FP8 block shares one scale across 16,384 weights; an NVFP4 group shares it across 16, so the grid follows a row's local range.

The engine's kernel is W4A16\index{W4A16}: it unpacks the 4-bit codes to 16-bit numbers in registers and multiplies on the ordinary tensor cores. On this board NVFP4 is a storage format, not a compute format. It saves bytes, and the arithmetic stays 16-bit (Chapter~\ref{ch:kernels}).

\section{The quantiser}\index{quantiser}\index{GPTQ}\index{clip search}

\code{tools/quant_nvfp4.py build} loads the BF16 model once (about 50~GB) and does everything in that process.

\begin{enumerate}
\item It runs a calibration text through the engine. \code{tools/calib_corpus.py} writes it from \code{bench/calib.txt}, Python standard-library modules, English Wikipedia and German Wikipedia: 149,682 tokens in the build that was published. Any document that shares a run of 12 words with a held-out file is dropped.
\item It records each projection's mean square input, and the full second moment of each input, $H = X^\top X$. At fp32 the second moments are 101~GB for 64 layers, so the build takes them in passes of about 13 layers (\code{--h-budget-gb 22}).
\item It quantises every projection with GPTQ on the NVFP4 grid. GPTQ rounds one input column at a time and pushes each column's rounding error onto the columns not yet rounded, weighted by the inverse of $H$. At the first column of each group of 16, a clip search picks the group's scale from the weights as they stand.
\end{enumerate}

The clip search tries five scale multipliers per group (1.00, 0.95, 0.90, 0.85 and 0.80). A smaller multiplier clips the group's largest weight and gives the rest a finer grid. It keeps the multiplier with the smallest error as the model sees it,
\[
\sum_k \textit{act\_ms}[k] \cdot \bigl(w_k - \text{dequant}(\text{quant}(w_k))\bigr)^2,
\]
where \textit{act\_ms}$[k]$ is the mean square of input channel $k$ over the calibration text. A channel the model never uses can take a coarse grid at no cost, and a channel it leans on cannot. GPTQ left a lower output error than the clip search alone on the calibration inputs, and a lower loss on held-out text, so the published set uses it. The format and the kernel do not change between the two; only the stored codes do.

The output head is quantised separately. \code{tools/quant_head.py build} writes it as e4m3 with one fp32 scale per vocabulary row (1.27~GB instead of 2.54~GB), searching the multipliers 1.0, 0.95 and 0.90 of each row's largest weight. One scale per row, because rows are tokens, and their sizes differ more than a block scale shared by 128 neighbouring ids can follow. The head decides what the engine writes, so it took the full gate. Its kernel returns fp32 logits at every row count.

\section{The gate}\label{sec:quality-gate}\index{quality gate}

\code{tools/quality_gate.py} reports three things, and a weight set ships only when all three hold.

\begin{enumerate}
\item Held-out loss, teacher-forced, as a change against a reference on the same tokens in the same process. With \code{--plan} the reference is a model loaded from another checkpoint, which is how the published set is scored against BF16 while it loads over the FP8 release as it does in serving. The gate files are \code{bench/heldout_code.txt}, a Python file the calibration does not contain, and \code{bench/heldout_prose.txt}, English prose on unrelated topics. The bar is +0.05 nats on the worse of the two.
\item Argmax agreement with the reference, overall and at positions where the reference is confident (its top choice ahead of the second by at least 1.0 in the logits). On a flat distribution over 248,320 tokens the top choice is decided by rounding, so the confident number is the one to watch.
\item A free generation of 900 greedy tokens on five prompts the calibration never saw, with repetition statistics.
\end{enumerate}

Test 3 matters most. Teacher-forced loss never lets an error grow, so it cannot see a model that drifts off its own path. On an earlier engine on this board, loss rated a broken configuration better than a healthy one.

\begin{figure}[htbp]
\centering
\begin{tikzpicture}
\begin{axis}[tandem, ybar, height=5.6cm, width=.9\linewidth, bar width=20pt,
    ymin=0, ymax=0.058, ylabel={held-out loss against BF16, nats/token},
    symbolic x coords={gp,gc,wp,wc8,wc10}, xtick=data,
    xticklabels={{gate prose\\2,048 tok},{gate code\\2,048 tok},{prose\\12k tok},{code\\8k tok},{more code\\10.5k tok}},
    x tick label style={align=center},
    yticklabel style={/pgf/number format/fixed, /pgf/number format/precision=2},
    scaled y ticks=false,
    nodes near coords, nodes near coords style={font=\sffamily\scriptsize,
      /pgf/number format/fixed, /pgf/number format/precision=4},
    enlarge x limits=0.15, grid=none]
  \addplot[fill=cobaltlight, draw=cobalt] coordinates
    {(gp,0.0221) (gc,0.0426) (wp,0.0174) (wc8,0.0184) (wc10,0.0272)};
  \draw[signal, thick, dashed] (rel axis cs:0,0.862) -- (rel axis cs:1,0.862)
    node[pos=0.98, above, font=\sffamily\scriptsize, text=signal, anchor=south east] {the bar: +0.05};
\end{axis}
\end{tikzpicture}
\caption[The published weights against BF16]{The published NVFP4 set against Qwen's BF16 release: the change in held-out loss on the two gate files and on three wider texts that neither the calibration nor anything else reads. \measured: the StairCut paper's quality section.}\label{fig:quality}
\end{figure}

The published set passes on both gate files: +0.0221 nats on prose and +0.0426 on code, inside the +0.05 bar (Figure~\ref{fig:quality}). Code costs more than prose. Re-run the gate before quantising anything further, and before trusting a weight set from elsewhere: a weight set passes or fails on its own measurement, not on its format.

\section{What stays at full precision}

The 48 recurrent states stay fp32. They are 151~MB, read and written once a step, about 1.5~ms of traffic. A recurrent state is a sum over the whole sequence, carried across every step, so its rounding error grows with every step; a weight's does not. The convolution weights, the decay and bias tensors of the recurrent layers, every norm and the drafters' own layers keep their shipped precision too.

\section{Profiles and bytes}

Three serving profiles ship (Table~\ref{tab:profiles-env}). The FP8 profile reads nothing but the vendor's files, so its quality is the checkpoint's. The balanced profile keeps the recurrent and attention layers on the vendor's weights and takes the published set's NVFP4 MLPs and FP8 head; the MLPs carry 70~\% of a verify's projection bytes. Neither profile's speed is measured on this release.

Bytes and accepted tokens partly substitute for each other. At one token a step every byte saved is time saved. On a round that commits several tokens the weight read is already spread over them, so the gain from fewer bytes shrinks as acceptance grows.

\section{Knobs}

\begin{table}[H]
\centering\small
\caption{Quantisation tools and the engine switches that load their output.}\label{tab:quant-knobs}
\begin{tabularx}{\linewidth}{@{}L{.48\linewidth}Y@{}}
\toprule
\hd{knob} & \hd{what} \\
\midrule
\code{--nvfp4 a,b,c} or \code{QWEN38_NVFP4} & the NVFP4 files to load, one per target; empty means the checkpoint's FP8 \\
\code{--fp8-head FILE}, \code{--fp8-head build}, \code{QWEN38_FP8_HEAD} & the e4m3 head; \code{build} makes it at load; empty means the BF16 head \\
\code{tools/quant_nvfp4.py build --methods gptq} & statistics and GPTQ from one load of the BF16 model \\
\code{tools/quant_head.py build --ratios} & the per-row scale search for the head \\
\code{tools/calib_corpus.py --out DIR --eval-chars} & the calibration text and the wider held-out texts \\
\code{tools/quant_sensitivity.py measure}, \code{mix --keep-gb} & per-piece sensitivity, and a mix that keeps the most sensitive pieces in FP8 \\
\code{tools/quality_gate.py --plan plan.json} & weight sets against a reference checkpoint \\
\bottomrule
\end{tabularx}
\end{table}

The gate files are one held-out file per domain, 2,048 tokens each, and the wider texts add 8k to 12k tokens per domain. No task benchmark suite has run against the quantised engine (Chapter~\ref{ch:measurement}, Limits).
